\section{Thesis organization}
\label{sec:Organization}
The remainder of this thesis is organized as follows.

Chapter~2 reviews the literature related to this research. It presents the fundamental concepts of cloud computing and hybrid cloud environments, followed by an overview of DevOps, DevSecOps, and SysSecOps. Existing security assessment approaches, risk evaluation methods, and current research on secure software delivery are also reviewed to establish the foundation and identify the limitations addressed by this research. Also, an AI-assisted code generation is introduced to increase code productivity.

Chapter~3 presents the proposed SysSecOps process for hybrid cloud environments. It describes the overall system architecture, the operational workflow, the unified security assessment mechanism, and the proposed risk evaluation process. This chapter also explains how continuous security is integrated throughout the software operational lifecycle.

Chapter~4 describes the implementation of the proposed SysSecOps system and evaluates its effectiveness. The implementation environment, prototype system, experimental scenarios, and evaluation methodology are presented, followed by the experimental results and performance analysis.

Chapter~5 discusses the findings obtained from the experimental evaluation. The proposed approach is analyzed with respect to the research objectives and existing studies, while its limitations and potential improvements are also discussed.

Finally, Chapter~6 concludes the thesis by summarizing the research, highlighting its main contributions, and suggesting directions for future work.